\section{Reedy structure on $\Lambda_P$}
\label{sec:interval-reedy}

In this section, we prove that $\Lambda_P$ admits a Reedy structure
with degree function given by the cardinality of intervals
(\Cref{thm:interval-cardinality-reedy-decomposition}).
The directional subalgebras $\Lambda_P^+$ and $\Lambda_P^-$ in this
decomposition are defined using quotient and inclusion maps associated with
relative downsets and relative upsets, respectively.

Each admissible-component basis morphism in
\Cref{prop:interval-hom-admissible-basis} factors through its image,
which is again an interval representation.  More precisely, let $R,S\in\Int(P)$ and
$U\in\Adm(R,S)$.  Then $U$ is a relative downset in $R$ and a relative
upset in $S$, and the image factorization is
\begin{equation}
\label{eq:admissible-component-factorization}
\rho_{R,S}^{U}
=
j_{U,S}\circ q_{R,U}
\colon
\I_R\twoheadrightarrow\I_U\hookrightarrow\I_S,
\end{equation}
where
$q_{R,U}:=\rho_{R,U}^{U}$ and
$j_{U,S}:=\rho_{U,S}^{U}$ are the corresponding quotient and inclusion.

More generally, the morphisms $q_{R,U}$ and $j_{U,S}$ are defined in the same way for
relative downsets $U\in\cD(R)$ and relative upsets
$U\in\cU(S)$, respectively.  In particular,
$q_{S,S}=j_{S,S}=\id_{\I_S}$.

Inside $E_P$, define the directional $\kk$-subspaces
\begin{equation}
\begin{aligned}
E_P^-
&=
\operatorname{span}_{\kk}
\{q_{R,U}\mid R\in\Int(P),\ U\in\cD(R)\},\\
E_P^+
&=
\operatorname{span}_{\kk}
\{j_{U,S}\mid S\in\Int(P),\ U\in\cU(S)\},\\
E_P^0
&=
\bigoplus_{S\in\Int(P)}\kk e_S.
\end{aligned}
\label{eq:interval-directional-subspaces}
\end{equation}
We then set
\begin{equation}
\Lambda_P^+
=
(E_P^-)^{\op},
\qquad
\Lambda_P^-
=
(E_P^+)^{\op},
\qand
\Lambda_P^0
=
E_P^0.
\label{eq:interval-Reedy-subspaces}
\end{equation}

The principal structural result of this section is the following.

\begin{theorem}
\label{thm:interval-cardinality-reedy-decomposition}
The triple
\[
(\Lambda_P,\Lambda_P^+,\Lambda_P^-)
\]
is a Reedy algebra with degree $\deg(S)=|S|$.  Its associated
quasi-hereditary order is given by reverse interval cardinality:
\begin{equation}
R\preceq_{\mathrm{qh}}S
\Longleftrightarrow
R=S
\text{ or }
|R|>|S|.
\label{eq:interval-cardinality-QH-order}
\end{equation}
\end{theorem}

We first verify the directional composition laws.

\begin{lemma}
\label{lem:interval-directional-composition}
Let $R,U,V\in\Int(P)$.
\begin{enumerate}[\rm (1)]
    \item If $U\in\cD(R)$ and $V\in\cD(U)$, then $V\in\cD(R)$ and $q_{U,V}\circ q_{R,U} = q_{R,V}$.

    \item If $V\in\cU(U)$ and $U\in\cU(R)$, then $V\in\cU(R)$ and $j_{U,R}\circ j_{V,U}=j_{V,R}$.
\end{enumerate}
\end{lemma}

\begin{proof}
Suppose that $U\in\cD(R)$ and $V\in\cD(U)$.  If $x\in R$ and
$x\leq v$ for some $v\in V$, then $x\in U$ because $U$ is a relative
downset in $R$, and hence $x\in V$ because $V$ is a relative downset in
$U$.  Thus $V\in\cD(R)$.  Both sides of the equality in {\rm (1)} are the identity on $V$ and
zero outside $V$, so they coincide.  The second assertion is dual.
\end{proof}

It follows that $E_P^-$ and $E_P^+$ are unital subalgebras of $E_P$
containing $E_P^0$.  Consequently, $\Lambda_P^+$ and $\Lambda_P^-$ are
unital subalgebras of $\Lambda_P$ containing $\Lambda_P^0$.  Moreover,
their diagonal blocks are
\[
e_S\Lambda_P^+e_S
=
\kk e_S
=
e_S\Lambda_P^-e_S.
\]
If $U\in\cD(R)$ and $U\neq R$, then $U\subsetneq R$, and
\[
q_{R,U}^{\op}\in e_R\Lambda_P^+e_U
\qand
|R|>|U|.
\]
Similarly, if $U\in\cU(S)$ and $U\neq S$, then $U\subsetneq S$, and
\[
j_{U,S}^{\op}\in e_U\Lambda_P^-e_S
\qand
|U|<|S|.
\]
Thus conditions {\rm (R1)} and {\rm (R2)} hold for
$\deg(S)=|S|$.  

\begin{proof}[Proof of \Cref{thm:interval-cardinality-reedy-decomposition}]
It remains to verify {\rm (R3)}.  Before passing to opposite algebras,
it is equivalent to show that multiplication induces an isomorphism
\begin{equation}
E_P^+
\otimes_{E_P^0}
E_P^-
\xrightarrow{\sim}
E_P.
\label{eq:interval-Reedy-multiplication-in-E}
\end{equation}

Fix $R,S\in\Int(P)$. 
The $(S,R)$-component of
\eqref{eq:interval-Reedy-multiplication-in-E} is
\begin{equation}
\bigoplus_{U\in\Int(P)}
e_SE_P^+e_U
\otimes_\kk
e_UE_P^-e_R
\longrightarrow
e_SE_Pe_R.
\label{eq:blockwise-interval-Reedy-multiplication}
\end{equation}
Its source has basis
\[
\{j_{U,S}\otimes q_{R,U}
\mid
U\in\cD(R)\cap\cU(S)\}.
\]
By \eqref{eq:admissible-components-as-relative-subsets}, the indexing
set is exactly $\Adm(R,S)$.  
For $U\in\Adm(R,S)$, 
\eqref{eq:admissible-component-factorization} gives 
$\rho_{R,S}^{U} =j_{U,S}\circ q_{R,U}$. 
These morphisms form a $\kk$-linear basis of
\[
e_SE_Pe_R
\cong
\Hom_P(\I_R,\I_S)
\]
by \Cref{prop:interval-hom-admissible-basis}.  Hence
\eqref{eq:blockwise-interval-Reedy-multiplication} is an isomorphism
for every $R,S\in\Int(P)$, and therefore
\eqref{eq:interval-Reedy-multiplication-in-E} is an isomorphism.
Passing to opposite algebras verifies {\rm (R3)}.

Thus $(\Lambda_P,\Lambda_P^+,\Lambda_P^-)$ is a Reedy algebra with
degree $\deg(S)=|S|$.  Its associated quasi-hereditary order is
\eqref{eq:interval-cardinality-QH-order} by \Cref{prop:Reedy-QH}.
\end{proof}

Throughout the remainder of the paper, we regard $\Lambda_P$ as equipped with the Reedy structure $(\Lambda_P,\Lambda_P^+,\Lambda_P^-)$ established above.


We first record the following composition law for the admissible-component basis. 

\begin{proposition}
\label{prop:admissible-basis-composition}
Let $Q,R,S\in\Int(P)$, and let
$D\in\Adm(Q,R)$ and $C\in\Adm(R,S)$.
Then every connected component $W$ of $C\cap D$ belongs to
$\Adm(Q,S)$, and
\begin{equation}
\rho_{R,S}^{C}\circ\rho_{Q,R}^{D}
=
\sum_{W\in\piZero(C\cap D)}
\rho_{Q,S}^{W},
\label{eq:admissible-basis-composition}
\end{equation}
where the empty sum is understood to be zero.
\end{proposition}

\begin{proof}
The subset $C\cap D$ is convex.  Hence each of its connected components
is an interval.  Let $W$ be such a component.  If $w\in W$ and $x\in Q$
satisfy $x\leq w$, then $x\in D$ because $D$ is a relative downset in
$Q$.  Since $x\in D\subseteq R$ and $w\in W\subseteq C$, the
relative-downset property of $C$ in $R$ gives $x\in C$.  Thus
$x\in C\cap D$.  Since $x$ and $w$ are comparable, they belong to the
same connected component, so $x\in W$.  Therefore $W\in\cD(Q)$.
Dually, $W\in\cU(S)$, and hence $W\in\Adm(Q,S)$ by
\eqref{eq:admissible-components-as-relative-subsets}.

The supports of the morphisms $\rho_{Q,S}^{W}$ are pairwise disjoint
and their union is $C\cap D$.  Thus
\eqref{eq:admissible-basis-composition} follows from
\eqref{eq:rho-definition}.
\end{proof}





Next, we consider the heredity chain determined by the Reedy degree and its
successive layers. For $0\leq t\leq |P|$, set
\begin{equation}
\varepsilon_t
:=
\sum_{\substack{T\in\Int(P)\\ |T|=t}} e_T,
\qquad
J_t
:=
\Lambda_P
(\varepsilon_0+\varepsilon_1+\cdots+\varepsilon_t)
\Lambda_P.
\label{eq:interval-cardinality-heredity-ideals}
\end{equation}
Since the intervals in $\Int(P)$ are nonempty, we have
$\varepsilon_0=0$ and hence $J_0=0$. By
\cite[Corollary~4.8]{CDK}, the sequence
\begin{equation}
0
=
J_0
\subseteq
J_1
\subseteq
\cdots
\subseteq
J_{|P|}
=
\Lambda_P
\label{eq:interval-cardinality-heredity-chain}
\end{equation}
is a heredity chain.
Since each $J_t$ is a two-sided ideal, we also regard this chain as a
filtration $J_\bullet$ of the regular
$\Lambda_P$--$\Lambda_P$-bimodule $\Lambda_P$.

By \eqref{eq:interval-cardinality-heredity-ideals} and
\Cref{prop:admissible-basis-composition}, we have
\begin{equation}
J_t
=
\operatorname{span}_{\kk}
\left\{
\bigl(\rho_{R,S}^{C}\bigr)^{\op}
\;\middle|\;
R,S,C\in\Int(P),\
C\in\Adm(R,S),\
|C|\leq t
\right\}.
\label{eq:interval-cardinality-admissible-basis-filtration}
\end{equation}
Since the admissible-component morphisms form a basis by
\Cref{prop:interval-hom-admissible-basis}, the quotient has the following basis description: 
\begin{equation}
J_t/J_{t-1}
=
\operatorname{span}_{\kk}
\left\{
(\rho_{R,S}^{C})^{\op}+J_{t-1}
\;\middle|\;
R,S,C\in\Int(P),\
C\in\Adm(R,S),\
|C|=t
\right\}.
\label{eq:interval-cardinality-admissible-basis-layer}
\end{equation}





\begin{remark}
\label{rem:reedy-versus-left-strongly-qh}
By the proof of \cite[Proposition~4.5]{AENY}, $\Lambda_P$ admits a
left strongly quasi-hereditary structure. 
The finiteness of
$\gldim\Lambda_P$ recorded in \Cref{prop:relative-auslander} follows from this structure.


By \Cref{thm:interval-cardinality-reedy-decomposition}, the Reedy
structure $(\Lambda_P,\Lambda_P^+,\Lambda_P^-)$ induces a
quasi-hereditary structure on $\Lambda_P$.  In general, this
quasi-hereditary structure is not left strongly quasi-hereditary, and
hence differs from the one considered in \cite{AENY}.
Indeed, let $P=\{a,b,c\}$ with $a<c$ and
$b<c$.  A direct calculation for the quasi-hereditary structure
induced by the Reedy structure gives the minimal projective resolution
\[
0
\longrightarrow
\PP(\{c\})
\longrightarrow
\PP(\{a,c\})\oplus\PP(\{b,c\})
\longrightarrow
\PP(P)
\longrightarrow
\DD(P)
\longrightarrow
0.
\]
Thus $\pd_{\Lambda_P}\DD(P)=2$, so this quasi-hereditary structure is
not left strongly quasi-hereditary.

The finiteness of $\gldim\Lambda_P$ also follows from
\Cref{thm:interval-cardinality-reedy-decomposition}, since
$\Lambda_P$ is quasi-hereditary.
\end{remark}
